% Resumen ejecutivo del Proyecto Integrador (plantilla institucional)
% Formato requerido: carta, una columna, 11 pt, márgenes de 2.5 cm, interlineado 1.15.
% Extensión: 3 páginas sin contar anexos. Referencias: APA 7.
% Compila con pdfLaTeX (opción predeterminada de Overleaf) + Biber.
\documentclass[letterpaper,11pt]{article}

\usepackage[utf8]{inputenc}
\usepackage[T1]{fontenc}
\usepackage{lmodern}
\usepackage[spanish,es-tabla,es-nodecimaldot]{babel}
\usepackage[margin=2.5cm]{geometry}
\usepackage{setspace}
\setstretch{1.15}
\usepackage{graphicx}
\usepackage{booktabs}
\usepackage{array}
\usepackage{float}
\usepackage[font=small,labelfont=bf]{caption}
\usepackage{csquotes}
\usepackage[hidelinks]{hyperref}

% Estilo de referencias: APA 7
\usepackage[style=apa,backend=biber]{biblatex}
\addbibresource{referencias.bib}

% --- Complementos del resumen ejecutivo ---
\usepackage{xcolor}
\setlength{\parindent}{1.5em}

% --- Encabezado: logos solo en la primera página ---
\usepackage{fancyhdr}
\geometry{headheight=52pt,headsep=0.3cm}
\renewcommand{\headrulewidth}{0pt}
\newcommand{\logo}[2]{%
  \IfFileExists{#1}%
    {\includegraphics[height=1.2cm]{#1}}%
    {\fbox{\parbox[c][1.1cm][c]{3cm}{\centering\sffamily\small #2}}}}
\fancypagestyle{primera}{%
  \fancyhf{}%
  \fancyhead[L]{\logo{logo.jpeg}{Logo Tec}}%
  \fancyhead[R]{\logo{sssa.png}{Logo empresa}}%
  \fancyfoot[C]{\thepage}}
\pagestyle{fancy}
\fancyhf{}
\fancyfoot[C]{\thepage}

% --- Recuadros de cifras (texto grande) ---
% \cifra{fondo}{color del texto}{valor}{etiqueta}{unidad}
\newcommand{\cifra}[5]{%
  \begingroup\setlength{\fboxsep}{5pt}%
  \colorbox{#1}{\parbox[t]{\dimexpr\linewidth-2\fboxsep\relax}{%
    \setstretch{1}\color{#2}{\Large\bfseries #3}\par
    #4%
    \ifx\relax#5\relax\else\par{\footnotesize #5}\fi}}%
  \endgroup}
\newcommand{\cifragris}[4]{\cifra{black!8}{black}{#1}{#2}{#3}}
\newcommand{\cifraacento}[4]{\cifra{blue!10}{blue!55!black}{#1}{#2}{#3}}

\begin{document}
\thispagestyle{primera}

% --- Título y datos ---
\noindent{\LARGE\bfseries Reconocimiento automático de actividades de manejo manual de cargas con pocas etiquetas\par}
\smallskip
\noindent{\large\textcolor{black!60}{Resumen ejecutivo}\par}
\medskip

\noindent\begingroup\setlength{\fboxsep}{7pt}%
\colorbox{black!8}{\parbox{\dimexpr\linewidth-2\fboxsep\relax}{\small\setstretch{1.1}%
\textcolor{black!60}{Integrantes:} Isaura Yutsil Flores Escamilla, Paola Alejandra León Guarneros, Dylan Ulises Quiroz Hernández\\
\textcolor{black!60}{Programa:} Proyecto Integrador · Maestría en Inteligencia Artificial Aplicada\\
\textcolor{black!60}{Asesor:} Gerardo Jesús Camacho González \quad \textcolor{black!60}{Fecha:} noviembre de 2026}}%
\endgroup\par\medskip

% --- Problema y solución ---
Levantar, sostener, transportar y colocar cargas a mano es una de las tareas industriales más asociadas con lesiones musculoesqueléticas. Para evaluar su riesgo, un área de ergonomía o de seguridad necesita saber \emph{qué} actividad realiza cada trabajador y \emph{cuándo}: no pesa lo mismo levantar una caja desde el piso que tomarla de una mesa. Hoy se obtiene marcando cada actividad a mano sobre video, un trabajo lento y costoso que limita la evaluación a muestras pequeñas de la jornada.

La solución propuesta es un sistema que \textbf{reconoce automáticamente siete actividades de manejo de cargas} a partir de \textbf{11 sensores de movimiento} colocados en el cuerpo, y entrega una línea de tiempo de actividades lista para el análisis ergonómico. Su diferencia frente a un modelo convencional es que aprende primero de horas de movimiento \emph{sin anotar}, con una técnica llamada codificación predictiva contrastiva (CPC) \parencite{vandenoord2018,haresamudram2021}, y después necesita muy pocos ejemplos etiquetados para adaptarse. Se validó con una base pública de 14 personas \parencite{bassani2021dataset}, siempre con personas que el modelo no había visto.

% --- Cifras de desempeño del modelo ---
\medskip
\noindent
\begin{minipage}[t]{0.235\linewidth}\cifragris{0.87}{F1 con personas nuevas}{media de 12 personas}{}\end{minipage}\hfill
\begin{minipage}[t]{0.235\linewidth}\cifragris{0.91 vs 0.90}{Frente al estado del arte}{misma persona de prueba}{}\end{minipage}\hfill
\begin{minipage}[t]{0.235\linewidth}\cifraacento{2.2$\times$}{Menos etiquetas}{con pocas anotaciones}{}\end{minipage}\hfill
\begin{minipage}[t]{0.235\linewidth}\cifragris{11 de 17}{Sensores necesarios}{sin perder desempeño}{}\end{minipage}
\medskip

% --- Desempeño ---
Con todas las etiquetas disponibles, el sistema alcanza un F1 macro de 0.87 en personas nuevas (1.0 sería perfecto), al nivel de un modelo supervisado bien ajustado (0.89) y del estudio publicado por los creadores de la base de datos: al replicar sus condiciones, nuestro modelo obtiene 0.91 en su persona de prueba frente a 0.90 de su red supervisada \parencite{bassani2025}. Su ventaja aparece donde más cuesta: \textbf{con 10 ejemplos etiquetados por actividad o menos, supera al modelo supervisado por hasta 0.11 de F1} en 11 o 12 de las 12 personas evaluadas, y para igualarlo el modelo supervisado necesita unas 2.2 veces más etiquetas. Además, casi nunca confunde una actividad a nivel de piso con una a nivel de mesa ($\le 0.1$\,\%), el error más costoso para un análisis ergonómico.

% --- Beneficios cuantificables (mensuales) ---
Para un área que hoy analiza 40 horas de grabación al mes, el escenario de referencia estima los siguientes beneficios mensuales:

\medskip
\noindent
\begin{minipage}[t]{0.485\linewidth}\cifragris{$-$140\,h}{Ahorro de tiempo de análisis}{horas/mes ($-$87.5\,\%)}{}\end{minipage}\hfill
\begin{minipage}[t]{0.485\linewidth}\cifragris{$-$\$35\,000}{Ahorro en costos operativos}{MXN/mes}{}\end{minipage}

\vspace{5pt}
\noindent
\begin{minipage}[t]{0.485\linewidth}\cifragris{+700\,\%}{Capacidad de procesamiento}{horas de grabación analizadas/mes, mismo personal}{}\end{minipage}\hfill
\begin{minipage}[t]{0.485\linewidth}\cifragris{$-$55\,\%}{Esfuerzo de anotación}{por cada adaptación a un área nueva (medido)}{}\end{minipage}

\smallskip
\noindent{\footnotesize\textcolor{black!60}{Cálculo y supuestos en el Anexo~A. Las cifras en pesos y horas son estimaciones que deben validarse con el sponsor; la reducción de anotación se midió en el proyecto.}}
\medskip

% --- Beneficios intangibles ---
\noindent\textbf{Mejora en la toma de decisiones.} La exposición se mide en turnos completos, no en unos minutos de video, y permite priorizar puestos con datos.\par
\smallskip
\noindent\textbf{Mejora en la calidad.} Todas las grabaciones se segmentan con el mismo criterio, sin variabilidad entre analistas.\par
\smallskip
\noindent\textbf{Facilidad o simplificación del trabajo.} El analista pasa de marcar video cuadro por cuadro a auditar una muestra.\par
\smallskip
\noindent\textbf{Innovación o diferenciación competitiva.} Aprender sin etiquetas y con 11 sensores en lugar de 17 acerca la evaluación ergonómica continua a la planta.\par
\medskip

% --- Inversión y costos (TCO a 12 meses) ---
El costo total de propiedad estimado para los primeros 12 meses es de \textbf{MXN 494\,000}: cinco meses para adaptar el modelo a sensores reales e implementar un piloto, y siete meses de operación (Anexo~B).

\medskip
\noindent
\begin{minipage}[t]{0.235\linewidth}\cifragris{\$183\,k}{Desarrollo}{MXN}{}\end{minipage}\hfill
\begin{minipage}[t]{0.235\linewidth}\cifragris{\$214\,k}{Implementación}{MXN}{}\end{minipage}\hfill
\begin{minipage}[t]{0.235\linewidth}\cifragris{\$97\,k}{Operación (7 meses)}{MXN}{}\end{minipage}\hfill
\begin{minipage}[t]{0.235\linewidth}\cifraacento{\$494\,k}{TCO 12 meses}{MXN}{}\end{minipage}
\medskip

Con el volumen de referencia, el primer año recupera cerca de la mitad de la inversión (MXN 245\,000 en siete meses de operación) y el retorno completo llega hacia el segundo año; con más horas analizadas, el retorno se acorta. \textbf{Recomendamos aprobar un piloto de cinco meses (MXN 397\,000 de desarrollo e implementación) en un área con manejo de cargas}, y condicionar la operación a escala a tres criterios: un F1 de al menos 0.85 en la auditoría del piloto, una latencia compatible con el uso en tiempo real y la confirmación del ahorro de horas de análisis.

% --- Riesgos y mitigación ---
Avanzar exige atender los riesgos que el propio proyecto identificó; la Tabla~\ref{tab:riesgos} resume los principales y cómo mitigarlos.

\begin{table}[H]
\centering
\footnotesize\setstretch{1}
\caption{Principales riesgos y medidas de mitigación.}
\label{tab:riesgos}
\begin{tabular}{@{}>{\raggedright\arraybackslash}p{0.44\linewidth}@{\hspace{0.04\linewidth}}>{\raggedright\arraybackslash}p{0.52\linewidth}@{}}
\toprule
\textbf{Riesgo} & \textbf{Mitigación}\\
\midrule
Base de 14 adultos jóvenes en laboratorio: el desempeño puede bajar en otra población o entorno. & Readaptar el modelo con grabaciones sin anotar del área y auditar una muestra antes de escalar.\\[1pt]
Los 11 sensores se simularon con un traje profesional; los comerciales tienen más ruido. & Probar el kit real en el piloto y medir cuánto ruido y desfase de calibración tolera el modelo.\\[1pt]
Validación final pendiente y, en su mayoría, con una sola semilla. & Completar la prueba con las dos personas reservadas y repetir con varias semillas.\\[1pt]
Colocar y levantar en la mesa tienen los errores más altos (F1 $\approx$ 0.82). & Usarlo para indicadores acumulados por turno, no para eventos aislados.\\[1pt]
Rechazo del personal o uso indebido de los datos. & Consentimiento informado, datos anónimos y uso sólo para mejorar puestos, no para evaluar personas.\\
\bottomrule
\end{tabular}
\end{table}

% --- Transferencia y continuidad ---
\section*{Transferencia y continuidad}
Se entregará al sponsor el repositorio del proyecto: código reproducible con su entorno, los modelos entrenados y los cuadernos de cada experimento con su diseño y resultados. Se incluyen el artículo técnico, una guía de instalación, el procedimiento para adaptar el sistema a un área nueva (grabar sin anotar y etiquetar de 10 a 50 ejemplos por actividad) y una propuesta de monitoreo que detecta descalibración y pérdida de desempeño sin etiquetas.

\section*{Declaración de uso de inteligencia artificial}
Se utilizó \textcite{anthropic2026claude} para redactar un borrador de este resumen a partir del artículo técnico del proyecto y para dar formato al documento. El equipo definió los supuestos de negocio, verificó las cifras y revisó la redacción final.
La responsabilidad final sobre el contenido entregado recae en los autores.

% --- Referencias (biblatex + Biber, solo lo citado) ---
\printbibliography[title={Referencias}]

% --- Anexos ---
\clearpage
\section*{Anexo A. Cálculo de los beneficios cuantificables}
Los beneficios se estiman para un escenario de referencia de un área que analiza manualmente grabaciones de manejo de cargas. Los supuestos S1--S4 son estimaciones del equipo que deben sustituirse por datos del sponsor; el indicador de anotación proviene de los experimentos del proyecto.

\begin{table}[H]
\centering
\small
\caption{Supuestos del escenario de referencia.}
\begin{tabular}{@{}l>{\raggedright\arraybackslash}p{9.4cm}r@{}}
\toprule
& Supuesto & Valor \\
\midrule
S1 & Grabación de manejo de cargas analizada al mes (p.\,ej., 20 puestos $\times$ 2~h) & 40 h \\
S2 & Tiempo de analista para segmentar manualmente 1~h de grabación & 4 h \\
S3 & Tiempo de analista para auditar 1~h de grabación segmentada por el sistema & 0.5 h \\
S4 & Costo por hora del analista, con prestaciones (MXN 40\,000/mes $\div$ 160~h) & MXN 250 \\
\bottomrule
\end{tabular}
\end{table}

\begin{table}[H]
\centering
\small
\caption{Beneficios mensuales estimados.}
\begin{tabular}{@{}>{\raggedright\arraybackslash}p{3.6cm}>{\raggedright\arraybackslash}p{2.6cm}>{\raggedright\arraybackslash}p{8.4cm}@{}}
\toprule
Indicador & Valor & Base de cálculo \\
\midrule
Ahorro de tiempo & 140 h/mes ($-$87.5\,\%) & S1 $\times$ (S2 $-$ S3) = 40 $\times$ (4 $-$ 0.5). Hoy: 160 h/mes; con el sistema: 20 h/mes de auditoría. \\
Ahorro en costos operativos & MXN 35\,000/mes & 140 h/mes $\times$ S4. \\
Capacidad de procesamiento & +700\,\% & Con las mismas 160 h/mes de analista se auditan 160 $\div$ 0.5 = 320 h de grabación, frente a 40 h hoy. \\
Esfuerzo de anotación al adaptar a un área nueva & $-$55\,\% por adaptación & Medido: para igualar lo que el sistema logra con 5 y 10 ejemplos por actividad, el modelo supervisado necesita unas 11 y 22 (2.2 veces más); 1 $-$ 1/2.2 $\approx$ 55\,\%. \\
\bottomrule
\end{tabular}
\end{table}

\clearpage
\section*{Anexo B. Costo total de propiedad (12 meses)}
Supuestos: ingeniero de aprendizaje automático con costo de MXN 60\,000/mes; sensores inerciales comerciales a USD 150 por unidad y un tipo de cambio de 18 MXN/USD; meses 1--3 de desarrollo, 4--5 de implementación y 6--12 de operación. Los precios son estimaciones de referencia que deben cotizarse.

\begin{table}[H]
\centering
\small
\caption{Desglose del costo total de propiedad por servicio, en MXN.}
\begin{tabular}{@{}l>{\raggedright\arraybackslash}p{4.3cm}>{\raggedright\arraybackslash}p{5.6cm}r@{}}
\toprule
Fase & Servicio & Base de cálculo & MXN \\
\midrule
Desarrollo & Adaptación del modelo a sensores reales & 1 ingeniero $\times$ 3 meses & 180\,000 \\
 & Anotación inicial para la adaptación & 50 ejemplos por actividad; 10~h de analista & 2\,500 \\
\cmidrule{2-4}
 & \multicolumn{2}{l}{\emph{Subtotal desarrollo}} & 182\,500 \\
\midrule
Implementación & Sensores & 2 kits de 11 sensores $\times$ USD 150 $\times$ 18 & 59\,400 \\
 & Equipo de cómputo local & 1 computadora portátil o servidor de borde & 25\,000 \\
 & Piloto e integración & 1 ingeniero $\times$ 2 meses & 120\,000 \\
 & Capacitación del personal & Sesiones para analistas y supervisores & 10\,000 \\
\cmidrule{2-4}
 & \multicolumn{2}{l}{\emph{Subtotal implementación}} & 214\,400 \\
\midrule
Operación & Mantenimiento y monitoreo & 20\,\% de un ingeniero $\times$ 7 meses & 84\,000 \\
 & Almacenamiento y respaldo & MXN 1\,000/mes $\times$ 7 meses & 7\,000 \\
 & Reposición de sensores & 10\,\% del costo de los kits & 6\,000 \\
\cmidrule{2-4}
 & \multicolumn{2}{l}{\emph{Subtotal operación}} & 97\,000 \\
\midrule
\multicolumn{3}{l}{\textbf{Costo total de propiedad a 12 meses}} & \textbf{493\,900} \\
\bottomrule
\end{tabular}
\end{table}

\end{document}
